% Ch12 WS3 -- mechanisms, collision model, catalysis + FRQ. Blocks 4-5.
\documentclass{apchem-worksheet}

\apcunitword{Chapter}
\apcunit{12}
\apcunittitle{Chemical Kinetics}
\apcdoctitle{Worksheet 3 \textbullet\ Mechanisms, Energy, Catalysis \& FRQ}
\apcsubtitle{Zumdahl \S12.5--12.7 \textbullet\ coefficients become exponents ONLY for elementary steps}

\begin{document}
\wsheader[When testing a mechanism, always run both checks: do the steps sum
to the overall equation, and does the rate-determining step reproduce the
experimental rate law?]

\problem[]{Write the rate law for each elementary step:}
\begin{parts}
  \item \ce{N2O -> N2 + O}: \answerblank[4cm]{$k[\ce{N2O}]$}
  \item \ce{NO + O3 -> NO2 + O2}:
        \answerblank[4.4cm]{$k[\ce{NO}][\ce{O3}]$}
  \item \ce{2 NOCl -> 2 NO + Cl2}:
        \answerblank[4cm]{$k[\ce{NOCl}]^2$}
  \item Why is it legitimate to use coefficients as exponents here, but not
        for an overall equation? \answerlines[3]{An elementary step
        describes an actual molecular collision, so its molecularity is its
        order. An overall equation summarizes net stoichiometry and says
        nothing about which collisions occur or in what order.}
\end{parts}

\problem[]{The overall reaction \ce{NO2 + CO -> NO + CO2} has the
experimental rate law rate $= k[\ce{NO2}]^2$. Consider the mechanism:
\begin{align*}
  \ce{NO2 + NO2 &-> NO3 + NO} &&\text{slow}\\
  \ce{NO3 + CO &-> NO2 + CO2} &&\text{fast}
\end{align*}}
\begin{parts}
  \item Show that the steps sum to the overall equation.
        \answerlines[2]{\ce{NO3} cancels, and one \ce{NO2} appears on both
        sides and cancels, leaving \ce{NO2 + CO -> NO + CO2}.}
  \item Identify the intermediate. \answerblank[2.4cm]{\ce{NO3}}
  \item Show that the mechanism predicts the observed rate law.
        \answerlines[2]{The slow step is bimolecular in \ce{NO2}, giving
        rate $= k[\ce{NO2}]^2$, which matches experiment.}
  \item Explain why \ce{CO} does not appear in the rate law even though it
        is a reactant. \answerlines[3]{\ce{CO} is consumed only in the fast
        step, which occurs \emph{after} the rate-determining step. Speeding
        up a step past the bottleneck cannot change the overall rate, so its
        concentration has no effect.}
\end{parts}

\problem[]{Explain Zumdahl's funnel analogy for the rate-determining step.
\answerlines[3]{Water poured rapidly through a funnel fills the container at
a rate set by the funnel opening, not by how fast you pour. Likewise the
overall reaction cannot proceed faster than its slowest step, regardless of
how fast the other steps are.}}

\problem[]{Collision model and energy:}
\begin{parts}
  \item State the two requirements for a collision to produce reaction.
        \answerlines[2]{Enough kinetic energy to meet or exceed the
        activation energy, and correct relative orientation of the reacting
        molecules.}
  \item A reaction is highly exothermic but extremely slow at room
        temperature. Explain how both can be true.
        \answerlines[3]{$\Delta H$ and $E_a$ are independent. A large
        negative $\Delta H$ says the products are much lower in energy, but
        a large $E_a$ means few collisions have enough energy to reach the
        transition state --- so the reaction is thermodynamically favourable
        yet kinetically slow.}
  \item Explain why a modest temperature increase can multiply the rate
        several-fold. \answerlines[3]{The fraction of molecules exceeding
        $E_a$ lies in the tail of the Maxwell--Boltzmann distribution, and
        that tail grows sharply with temperature. The Arrhenius factor
        $e^{-E_a/RT}$ makes the dependence exponential rather than linear.}
\end{parts}

\problem[]{Catalysis:}
\begin{parts}
  \item How does a catalyst increase a reaction rate?
        \answerlines[2]{It provides an alternative mechanism with a lower
        activation energy, so a larger fraction of collisions succeed.}
  \item Name three things a catalyst does NOT change.
        \answerblank[7.1cm]{$\Delta H$, equilibrium position, products}
  \item In this mechanism, classify each species:
        \ce{Cl + O3 -> ClO + O2}; \ce{ClO + O -> Cl + O2}.
        \answerlines[2]{\ce{Cl} is a catalyst --- consumed first, then
        regenerated. \ce{ClO} is an intermediate --- produced first, then
        consumed.}
  \item Give the rule that distinguishes them.
        \answerlines[2]{Order of appearance: a catalyst appears first as a
        reactant and is regenerated; an intermediate appears first as a
        product and is later consumed.}
\end{parts}

\problem[]{\textbf{FRQ (10 points).} The gas-phase reaction
\ce{2 NO(g) + O2(g) -> 2 NO2(g)} was studied at \SI{25}{\celsius}.}

\renewcommand{\arraystretch}{1.25}
\begin{center}
\begin{tabular}{@{}cccc@{}}
\toprule
Trial & $[\ce{NO}]$ (M) & $[\ce{O2}]$ (M) & initial rate (\si{\Molar\per\second})\\
\midrule
1 & 0.010 & 0.010 & $2.5\times10^{-5}$\\
2 & 0.020 & 0.010 & $1.0\times10^{-4}$\\
3 & 0.010 & 0.020 & $5.0\times10^{-5}$\\
\bottomrule
\end{tabular}
\end{center}

\begin{parts}
  \item Determine the order with respect to each reactant, showing the
        trials compared. \answerlines[3]{Trials 1 and 2 ($[\ce{O2}]$ fixed):
        [NO] doubled, rate $\times 4$, so second order in NO.
        Trials 1 and 3 ([NO] fixed): $[\ce{O2}]$ doubled, rate doubled, so
        first order in \ce{O2}.}
  \item Write the rate law and calculate $k$ with units.
        \workspace[2.2cm]
        \keyonly{\begin{apcanswer}rate $= k[\ce{NO}]^2[\ce{O2}]$;
        $k = 2.5\times10^{-5}/[(0.010)^2(0.010)] =
        \SI{25}{\per\Molar\squared\per\second}$\end{apcanswer}}
  \item Predict the initial rate when $[\ce{NO}] = \SI{0.030}{\Molar}$ and
        $[\ce{O2}] = \SI{0.015}{\Molar}$. \workspace[2cm]
        \keyonly{\begin{apcanswer}rate $= 25(0.030)^2(0.015) =
        \SI{3.4e-4}{\Molar\per\second}$\end{apcanswer}}
  \item The following mechanism is proposed:
        \begin{align*}
          \ce{NO + NO &-> N2O2} &&\text{fast equilibrium}\\
          \ce{N2O2 + O2 &-> 2 NO2} &&\text{slow}
        \end{align*}
        Identify the intermediate and verify that the steps sum to the
        overall equation. \answerlines[2]{\ce{N2O2} is the intermediate; it
        cancels, leaving \ce{2 NO + O2 -> 2 NO2}.}
  \item Show that this mechanism is consistent with your rate law from (b).
        \answerlines[3]{The slow step gives rate $=
        k[\ce{N2O2}][\ce{O2}]$, but an intermediate cannot appear in a rate
        law. The fast pre-equilibrium makes $[\ce{N2O2}]$ proportional to
        $[\ce{NO}]^2$, so substituting gives rate $\propto
        [\ce{NO}]^2[\ce{O2}]$ --- matching the experimental law.}
\end{parts}

\begin{apcnote}[Rubric --- score yourself, 10 points]
\textbf{(a) 3 pts} 1 per order, plus 1 for naming the trials compared and
what was held constant --- orders without that reasoning earn half
\textbullet\ \textbf{(b) 2 pts} 1 rate law, 1 value of $k$ WITH units
\textbullet\ \textbf{(c) 1 pt} $3.4\times10^{-4}$ \si{\Molar\per\second}
\textbullet\ \textbf{(d) 2 pts} 1 intermediate, 1 steps summed
\textbullet\ \textbf{(e) 2 pts} 1 states that an intermediate cannot appear
in a rate law, 1 performs the pre-equilibrium substitution.
\end{apcnote}

\begin{spiralstrip}{Chapter 12 \textbullet\ blocks 1--3}
  \item $t_{1/2}$ for first order with $k = \SI{0.0347}{\per\second}$:
        \answerblank[2.8cm]{\SI{20.0}{\second}}
  \item Which plot is linear for zero order?
        \answerblank[2.6cm]{$[\ce{A}]$ vs $t$}
  \item Fraction left after 4 half-lives: \answerblank[1.8cm]{1/16}
  \item Rate law for the elementary step \ce{A + A -> B}:
        \answerblank[2.4cm]{$k[\ce{A}]^2$}
  \item $E_a$ affects the rate or the yield? \answerblank[2cm]{the rate}
\end{spiralstrip}

\end{document}
